\documentclass[rapport]{subfiles}
\renewcommand{\author}{Thibauld, Benjamin, Amritesh}
\begin{document}
	\section{Définition du problème statistique}
	\subsection{Population}
	L’ensemble des individus sur lequel porte notre étude est la classe de CPI 1.
	\subsection{Échantillon}
	Notre échantillon est constitué de ... étudiant.
	\subsection{Variable aléatoire}
	La variable aléatoire étudié est le nombre d'heure de travail.
	\section{Collecte des données}
	\begin{table}[H]
		\centering
		\caption{Tableau des résultats}
		\label{tab:1}
		\begin{tabular}{|c|p{3cm}|p{3cm}|p{3cm}|}
			\hline
			& travail en jour de semaine  & travail en jour weekend & travail en vacances \\
			\hline
			$[0;1[$ &  &  &  \\
			\hline
			$[1;2[$ &  &  &  \\
			\hline
			$[2;3[$ &  &  &  \\
			\hline
			$[3;4[$ &  &  &  \\
			\hline
			$[4;5[$ &  &  &  \\
			\hline
			$[5;6[$ &  &  &  \\
			\hline
			$[6;7[$ &  &  &  \\
			\hline
			+ de 7 &  &  &  \\
			\hline
		\end{tabular}
	\end{table}

	\section{Analyse statistique avec Python}
	\section{Représentations graphiques}
	\subsection{Histogramme}
	\subsection{Boîte à moustaches}
	\section{Modélisation probabiliste}
	\section{Étude de la loi avec Python}
	\section{Interprétation}
\end{document}